\textbf{Phần 1. Trắc nghiệm lựa chọn}
\setcounterpageref{ex}{0}

\begin{ex} %[5D7N40]
Để tìm tỉ số phần trăm của hai số $12$ và $25$, ta thực hiện
\choice
{$12 \times 25 : 100$}
{$25 : 12 \times 100\%$}
{$(12 + 25) : 100\%$}
{\True $12 : 25 \times 100\%$}
\loigiai{
Để tìm tỉ số phần trăm của $12$ và $25$, ta lấy $12 : 25 = 0{,}48 = 48\%$.\\
\textbf{Đáp án của câu hỏi là D}
}
\end{ex}

\begin{ex} %[5D7N40]
Tỉ số phần trăm của hai số $19$ và $30$ làm tròn đến hàng phần mười là
\choice
{\True $63{,}3\%$}
{$63\%$}
{$63{,}33\%$}
{$63{,}4\%$}
\loigiai{
Ta có $19 : 30 \approx 0{,}6333... = 63{,}33...\% \approx 63{,}3\%$.\\
\textbf{Đáp án của câu hỏi là A}
}
\end{ex}

\begin{ex} %[5D7H40]
Một lớp học có $40$ học sinh, trong đó có $26$ bạn nữ. Tỉ số phần trăm của số học sinh nữ so với học sinh cả lớp là
\choice
{$60\%$}
{$52\%$}
{$35\%$}
{\True $65\%$}
\loigiai{
Tỉ số phần trăm học sinh nữ là: $26 : 40 = 0{,}65 = 65\%$.\\
\textbf{Đáp án của câu hỏi là D}
}
\end{ex}

\begin{ex} %[5D7H40]
Một đàn gà có $150$ con, trong đó có $60$ con gà trống. Tỉ số phần trăm của số gà mái so với tổng số gà là
\choice
{$40\%$}
{\True $60\%$}
{$50\%$}
{$75\%$}
\loigiai{
Số gà mái là: $150 - 60 = 90$ (con).\\
Tỉ số phần trăm số gà mái: $90 : 150 = 0{,}6 = 60\%$.\\
\textbf{Đáp án của câu hỏi là B}
}
\end{ex}

\begin{ex} %[5D7H40]
Theo kế hoạch, tổ công nhân phải may $500$ chiếc áo, tổ đã may được $560$ chiếc áo. Tổ công nhân đã thực hiện được bao nhiêu phần trăm kế hoạch?
\choice
{$110\%$}
{$120\%$}
{\True $112\%$}
{$125\%$}
\loigiai{
Tỉ số phần trăm thực hiện được là: $560 : 500 = 1{,}12 = 112\%$.\\
\textbf{Đáp án của câu hỏi là C}
}
\end{ex}

\begin{ex} %[5D7H40]
Một xưởng may sản xuất được $480$ sản phẩm đạt chuẩn trên tổng số $500$ sản phẩm đã may. Tỉ số phần trăm sản phẩm đạt chuẩn là
\choice
{\True $96\%$}
{$94\%$}
{$95\%$}
{$98\%$}
\loigiai{
Tỉ số phần trăm sản phẩm đạt chuẩn là: $480 : 500 = 0{,}96 = 96\%$.\\
\textbf{Đáp án của câu hỏi là A}
}
\end{ex}

\begin{ex} %[5D7V40]
Một người mua một chiếc xe đạp với giá $1\,600\,000$ đồng rồi bán lại với giá $1\,840\,000$ đồng. Người đó đã lãi được bao nhiêu phần trăm so với tiền vốn?
\choice
{$12\%$}
{\True $15\%$}
{$18\%$}
{$24\%$}
\loigiai{
Số tiền lãi là: $1\,840\,000 - 1\,600\,000 = 240\,000$ (đồng).\\
Tỉ số phần trăm tiền lãi so với vốn là: $240\,000 : 1\,600\,000 = 0{,}15 = 15\%$.\\
\textbf{Đáp án của câu hỏi là B}
}
\end{ex}

\begin{ex} %[5D7V40]
Một mảnh đất hình chữ nhật có chiều dài $25\text{ m}$ và chiều rộng $16\text{ m}$. Người ta dành $120\text{ m}^2$ diện tích đất để đào ao. Diện tích ao chiếm bao nhiêu phần trăm diện tích mảnh đất?
\choice
{$25\%$}
{$35\%$}
{$40\%$}
{\True $30\%$}
\loigiai{
Diện tích mảnh đất: $25 \times 16 = 400\text{ (m}^2\text{)}$.\\
Tỉ số phần trăm diện tích ao: $120 : 400 = 0{,}3 = 30\%$.\\
\textbf{Đáp án của câu hỏi là D}
}
\end{ex}

\begin{ex} %[5D7V40]
Khối 5 của một trường tiểu học có $200$ học sinh. Qua khảo sát có $130$ bạn thích bơi, $50$ bạn thích bóng rổ, còn lại thích cờ vua. Tỉ số phần trăm học sinh thích cờ vua là
\choice
{$5\%$}
{$15\%$}
{\True $10\%$}
{$20\%$}
\loigiai{
Số học sinh thích cờ vua là: $200 - (130 + 50) = 20$ (bạn).\\
Tỉ số phần trăm: $20 : 200 = 0{,}1 = 10\%$.\\
\textbf{Đáp án của câu hỏi là C}
}
\end{ex}

\begin{ex} %[5D7V40]
Một sản phẩm có giá gốc là $80\,000$ đồng, cửa hàng bán với giá $100\,000$ đồng. Tỉ số phần trăm giữa tiền lãi và tiền bán là
\choice
{$25\%$}
{\True $20\%$}
{$15\%$}
{$30\%$}
\loigiai{
Tiền lãi: $100\,000 - 80\,000 = 20\,000$ (đồng).\\
Tỉ số phần trăm giữa tiền lãi và tiền bán: $20\,000 : 100\,000 = 20\%$.\\
\textbf{Đáp án của câu hỏi là B}
}
\end{ex}

\begin{ex} %[5D7V40]
Nước biển chứa $3{,}5\%$ muối. Để lấy được $140\text{ kg}$ muối thì cần làm bay hơi bao nhiêu tấn nước biển?
\choice
{\True $4$ tấn}
{$2$ tấn}
{$3$ tấn}
{$5$ tấn}
\loigiai{
Lượng nước biển cần là: $140 : 3{,}5 \times 100 = 4\,000\text{ (kg)} = 4\text{ tấn}$.\\
\textbf{Đáp án của câu hỏi là A}
}
\end{ex}

\begin{ex} %[5D7C40]
Một cửa hàng nhập về một lô hàng với giá $100$ triệu đồng. Đợt 1 bán được $\dfrac{3}{5}$ số hàng với mức lãi $20\%$. Đợt 2 bán hết số hàng còn lại bị lỗ $10\%$. Sau khi bán hết lô hàng, cửa hàng đó lãi được bao nhiêu phần trăm so với vốn ban đầu?
\choice
{$10\%$}
{$6\%$}
{\True $8\%$}
{$12\%$}
\loigiai{
Vốn đợt 1: $100 \times \dfrac{3}{5} = 60$ triệu. Tiền lãi đợt 1: $60 \times 20\% = 12$ triệu.\\
Vốn đợt 2: $100 - 60 = 40$ triệu. Tiền lỗ đợt 2: $40 \times 10\% = 4$ triệu.\\
Tổng lãi thực tế: $12 - 4 = 8$ triệu.\\
Tỉ số phần trăm tiền lãi so với vốn ban đầu: $8 : 100 = 8\%$.\\
\textbf{Đáp án của câu hỏi là C}
}
\end{ex}

\textbf{Phần 2. Trắc nghiệm đúng sai}
\setcounterpageref{ex}{0}

\begin{ex} %[5D7H40]
Một tổ sản xuất theo kế hoạch phải làm $400$ chi tiết máy, nhưng thực tế tổ đã làm được $460$ chi tiết máy. Xét tính đúng sai:
\choiceTF
{\True Tổ đó đã thực hiện được $115\%$ so với kế hoạch.}
{Tổ đó đã vượt mức kế hoạch là $60\%$.}
{\True Tổ đó đã vượt mức kế hoạch là $15\%$.}
{Nếu tổ chỉ làm được $360$ chi tiết thì tổ đã đạt $80\%$ kế hoạch.}
\loigiai{
\begin{itemize}
    \item Thực tế so với kế hoạch: $460 : 400 = 1{,}15 = 115\%$. Đúng.
    \item Mức vượt: $115\% - 100\% = 15\% \neq 60\%$. Sai.
    \item Mức vượt là $15\%$. Đúng.
    \item Nếu làm 360: $360 : 400 = 90\% \neq 80\%$. Sai.
\end{itemize}
\begin{itemize}
    \item a) \textbf{Đúng}.
    \item b) \textbf{Sai}.
    \item c) \textbf{Đúng}.
    \item d) \textbf{Sai}.
\end{itemize}
}
\end{ex}

\begin{ex} %[5D7V40]
Một bình nước muối chứa $400\text{ g}$ dung dịch có $20\text{ g}$ muối. Xét tính đúng sai:
\choiceTF
{\True Tỉ số phần trăm muối trong dung dịch là $5\%$.}
{Nếu đổ thêm $100\text{ g}$ nước lọc vào bình thì tỉ số phần trăm muối tăng lên.}
{\True Sau khi đổ thêm $100\text{ g}$ nước lọc, dung dịch mới chứa $4\%$ muối.}
{Nếu cho thêm $5\text{ g}$ muối vào dung dịch ban đầu thì tỉ số phần trăm muối là $6\%$.}
\loigiai{
\begin{itemize}
    \item Tỉ lệ muối ban đầu: $20 : 400 = 5\%$. Đúng.
    \item Thêm nước lọc làm dung dịch loãng ra nên nồng độ giảm đi, không thể tăng lên. Sai.
    \item Thêm 100g nước lọc: khối lượng dung dịch là $500\text{ g}$, tỉ lệ muối: $20 : 500 = 4\%$. Đúng.
    \item Thêm 5g muối: muối là $25\text{ g}$, dung dịch là $405\text{ g}$, tỉ lệ $25 : 405 \approx 6{,}17\% \neq 6\%$. Sai.
\end{itemize}
\begin{itemize}
    \item a) \textbf{Đúng}.
    \item b) \textbf{Sai}.
    \item c) \textbf{Đúng}.
    \item d) \textbf{Sai}.
\end{itemize}
}
\end{ex}

\textbf{Phần 3. Trả lời ngắn}
\setcounterpageref{ex}{0}

\begin{ex} %[5D7H40]
Tìm tỉ số phần trăm của hai số $33$ và $75$. (Chỉ ghi số).
\shortans{44}
\loigiai{
\textbf{Đáp số:} $44$\\
$33 : 75 = 0{,}44 = 44\%$.
}
\end{ex}

\begin{ex} %[5D7H40]
Một trang trại có $800$ con vịt, trong đó có $360$ con vịt đực. Hỏi số vịt cái chiếm bao nhiêu phần trăm tổng số vịt? (Chỉ ghi số).
\shortans{55}
\loigiai{
\textbf{Đáp số:} $55$\\
Số vịt cái: $800 - 360 = 440$. Tỉ số phần trăm: $440 : 800 = 55\%$.
}
\end{ex}

\begin{ex} %[5D7V40]
Một cửa hàng nhập về một món đồ với giá $500\,000$ đồng và bán với giá $625\,000$ đồng. Hỏi cửa hàng đã lãi bao nhiêu phần trăm so với tiền vốn? (Chỉ ghi số).
\shortans{25}
\loigiai{
\textbf{Đáp số:} $25$\\
Tiền lãi: $125\,000$. Tỉ số phần trăm lãi: $125\,000 : 500\,000 = 25\%$.
}
\end{ex}

\begin{ex} %[5D7V40]
Cuối năm học, lớp 5C có $36$ học sinh hoàn thành xuất sắc nhiệm vụ học tập trên tổng số $45$ học sinh cả lớp. Tính tỉ số phần trăm học sinh hoàn thành xuất sắc nhiệm vụ học tập của lớp 5C. (Chỉ ghi số).
\shortans{80}
\loigiai{
\textbf{Đáp số:} $80$\\
$36 : 45 = 0{,}8 = 80\%$.
}
\end{ex}

\begin{ex} %[5D7V40]
Một người gửi tiết kiệm $50\,000\,000$ đồng, sau một năm nhận được cả tiền gốc và tiền lãi là $53\,500\,000$ đồng. Hỏi lãi suất gửi tiết kiệm một năm là bao nhiêu phần trăm? (Chỉ ghi số thập phân có dấu phẩy nếu có).
\shortans{7}
\loigiai{
\textbf{Đáp số:} $7$\\
Tiền lãi: $3\,500\,000$ đồng. Lãi suất: $3\,500\,000 : 50\,000\,000 = 7\%$.
}
\end{ex}

\begin{ex} %[5D7C40]
Bán một chiếc quạt điện với giá $420\,000$ đồng thì được lãi $20\%$ so với giá vốn. Hỏi tiền vốn của chiếc quạt điện đó là bao nhiêu nghìn đồng? (Ghi kết quả dưới dạng số nghìn đồng, ví dụ: 350).
\shortans{350}
\loigiai{
\textbf{Đáp số:} $350$\\
Giá bán chiếm $100\% + 20\% = 120\%$ tiền vốn.\\
Tiền vốn: $420\,000 : 120 \times 100 = 350\,000$ (đồng) $= 350$ nghìn đồng.
}
\end{ex}
